\subsection{2. Lokale Algebra und ehrlicher
Interventionstest}\label{lokale-algebra-und-ehrlicher-interventionstest}

Die Senderoperation ist

\[\adjustbox{max width=\linewidth}{$\displaystyle 
Z_4=(-1)^{n_4}=I-2n_4,\qquad
\mathcal E_4(\rho)=Z_4\rho Z_4^\dagger.
$}\]

Sie hat genau einen unitären Krausoperator und ist daher vollständig
positiv und spurerhaltend. Sie ist fermionparitätsgerade und erhält N.
Auf der ganzen Fock-Algebra kommutiert sie mit allen disjunkten
Modenoperatoren; insbesondere \([Z_4,n_5]=0\). Ihre Einschränkung auf
den ausgewählten Block wechselt allein das Vorzeichen des Blattes
(4,57). Diese Einschränkung wurde gegen den wirklichen Operator
\((-1)^{n_4}\) auf allen 2076 N=2-Basiszuständen geprüft.

Der Empfänger liest \(B=n_5\) aus. Im gewählten invarianten Block ist
dies der Projektor auf Blatt (5,56). Die Operation enthält keinen
Verweis auf Mode 5. Zu ihrem Zeitpunkt bleibt dessen gesamte
Besetzungsstatistik für jeden Zustand identisch, weil
\(Z_4^\dagger n_5 Z_4=n_5\).

Mit \(\rho_T=U_T|p_0\rangle\langle p_0|U_T^\dagger\) ist die gemessene
Differenz genau

\[\adjustbox{max width=\linewidth}{$\displaystyle 
\delta_{4\to5}(T)=
\operatorname{tr}\!\left[n_5U_T\mathcal E_4(\rho_T)U_T^\dagger\right]
-\operatorname{tr}\!\left[n_5U_T\rho_TU_T^\dagger\right].
$}\]

Das ist ein Vergleich von unbedingten Wahrscheinlichkeiten nach einem
kontrollierten lokalen Eingriff. Eine Zweipunktkorrelation wird dafür
nicht als Ersatz benutzt. Die vor dem Impuls erzeugte Kohärenz wird
ausdrücklich durch die native Vorentwicklung präpariert. Ein
Phasenimpuls direkt auf dem anfänglichen Produktzustand hätte nur ein
globales Vorzeichen und überhaupt keinen Effekt; diese Negativkontrolle
wird ebenfalls geprüft.

\subsection{3. Exakter Test mit gleicher
Endzusammensetzung}\label{exakter-test-mit-gleicher-endzusammensetzung}

Setze für \(\Delta>0\) und reelles \(g\ne0\)

\[\adjustbox{max width=\linewidth}{$\displaystyle 
\Omega=\sqrt{\Delta^2/4+8g^2},\quad
T=\frac\pi\Omega,\quad d=\frac\Delta{2\Omega},\quad
z=-e^{-i\pi d},\quad
x=|z-1|^2=4\cos^2\frac{\pi d}{2}.
$}\]

Nach T wirkt der helle Zweizustandsblock als z mal Identität; die
dunklen Paarzustände bleiben unverändert. Es gibt dann für \textbf{jede}
anfängliche Paarkombination keine Bosonamplitude. Auf der Paarseite ist

\[\adjustbox{max width=\linewidth}{$\displaystyle 
U_T^{(f)}=I_8+(z-1)|S\rangle\langle S|.
$}\]

Beide Arme beginnen in demselben vollständig angegebenen Produktzustand
\(|p_0\rangle=f_4^\dagger f_{57}^\dagger|0\rangle\), alle anderen Moden
leer. Sie unterscheiden sich nur durch den mittleren Phasenimpuls:

\[\adjustbox{max width=\linewidth}{$\displaystyle 
\begin{array}{lll}
\text{unberührt:}&U_T\;I\;U_T|p_0\rangle,&\text{danach }n_5,\\
\text{Eingriff:}&U_T\;Z_4\;U_T|p_0\rangle,&\text{danach }n_5.
\end{array}
$}\]

Für jedes andere Blatt q sind die exakten Endamplituden

\[\adjustbox{max width=\linewidth}{$\displaystyle 
\langle p_q|U_T^2|p_0\rangle
=s_0s_q\frac{z^2-1}{8},\qquad
\langle p_q|U_TZ_4U_T|p_0\rangle
=s_0s_q\frac{3(z-1)^2}{32}.
$}\]

Damit lautet insbesondere die unbedingte Empfängerwahrscheinlichkeit

\[\adjustbox{max width=\linewidth}{$\displaystyle 
p_{\rm frei}(n_5=1)=\frac{x(4-x)}{64},\qquad
p_{\rm Impuls}(n_5=1)=\frac{9x^2}{1024},
$}\]

und die kausale Differenz ist

\[\adjustbox{max width=\linewidth}{$\displaystyle 
\boxed{\delta_{4\to5}(T)=\frac{x(25x-64)}{1024}}.
$}\]

Für \(0<g^2\le3\Delta^2/32\) gilt \(1/2\le d<1\), also \(0<x\le2\) und
folglich \textbf{strikt \(\delta_{4\to5}<0\)}. Das ist eine exakte
endliche Zeit, kein unkontrollierter Schluss aus einer numerisch kleinen
Antwort oder einer asymptotischen Reihe. Außerhalb dieses Intervalls
gilt weiterhin die exakte Formel, kann aber auch das Vorzeichen wechseln
oder an x=64/25 verschwinden.

Am bisherigen Prüfpunkt \(\Delta=1,g=1/20\) ist \(d=5/(3\sqrt3)\) und
\(T\approx6.0459978807807\):

{\def\LTcaptype{none} % do not increment counter
\begin{longtable}[]{@{}lr@{}}
\toprule\noalign{}
Größe & Wert \\
\midrule\noalign{}
\endhead
\bottomrule\noalign{}
\endlastfoot
Ohne mittleren Impuls & 0.00087491599417248 \\
Mit mittlerem Impuls & 0.0000017344871305812 \\
Differenz & −0.00087318150704190 \\
Endbesetzung in beiden Armen & Nf=2, Nb=0 exakt \\
\end{longtable}
}

Die Wellenfunktion kann zwischen den Ablesezeiten einen Boson enthalten.
Das natürliche H vermittelt dadurch die Umverteilung von einem
Fermionenpaar in andere. Die Endmessung benötigt keine Selektion auf
„kein Boson``: Dieser Zustand gilt in beiden Armen mit
Wahrscheinlichkeit eins.

\subsection{4. Zwei ergänzende
Kontrollen}\label{zwei-erguxe4nzende-kontrollen}

\subsubsection{Boson als einfachere
Referenz}\label{boson-als-einfachere-referenz}

Ausgehend vom Produktzustand \(b_0^\dagger|0\rangle\) wähle die beiden
Verzögerungen \(\tau=\pi/(2\Omega)\). Der unberührte Arm kehrt exakt zum
Boson zurück, also \(p_{\rm frei}(n_5=1)=0\). Ein mittleres Z4 erzeugt

\[\adjustbox{max width=\linewidth}{$\displaystyle 
p_{\rm Impuls}(n_5=1)
=\frac{(1-d^2)\left[1+d^2-2d\sin(\pi d/2)\right]}{128}
\ge\frac{(1-d^2)(1-d)^2}{128}>0.
$}\]

Dies gilt für jedes \(g\ne0\) und \(\Delta>0\). Auch dieser Test ist
ladungserhaltend und unbedingt. Er hat allerdings unterschiedliche
Fermion/Boson-Zusammensetzungen am Ende; für die stärkere Aussage über
reine Paarumverteilung wird deshalb der Test aus Abschnitt 3 benutzt.

\subsubsection{Lokales Auffüllen und kurze
Zeiten}\label{lokales-auffuxfcllen-und-kurze-zeiten}

Auf der Referenz \(f_{57}^\dagger|0\rangle\) füllt das lokale
CPTP-Instrument mit Krausoperatoren \(n_4,f_4^\dagger\) Mode 4. Ohne
Eingriff ist der Einfermionzustand stationär. Mit Eingriff entsteht
\(|p_0\rangle\), woraus

\[\adjustbox{max width=\linewidth}{$\displaystyle 
\delta n_5(t)=\frac{|a(t)-1|^2}{64},\qquad
a(t)=e^{-i\Delta t/2}\left[\cos\Omega t+i d\sin\Omega t\right].
$}\]

Die führenden Beiträge sind \(\delta n_5(t)=g^4t^4/4+O(t^6)\) und
\(\delta N_{b_0}(t)=g^2t^2+O(t^4)\). Der Pfad in ein anderes
Fermionenpaar hat somit zwei native H-Anwendungen in der Amplitude; die
erste Umwandlung zum Boson hat eine.

Dieses Auffüllen tauscht eine Ladungseinheit mit einem angenommenen
Reservoir aus. Es wird nicht als zahlenerhaltende Operation ausgegeben
und ist für den Hauptnachweis nicht erforderlich.

\subsection{5. Prüfung und unveränderte offene
Grenze}\label{pruxfcfung-und-unveruxe4nderte-offene-grenze}

\codeword{verify\_causal.py} benutzt explizite Fehlerbedingungen, keine
abschaltbaren Assertions. Normaler Python-Lauf und
\codeword{python -OO}, jeweils mit Warnungen als Fehler, erzeugen
byte-identische JSON-Ergebnisse. Der optimierte Lauf wurde aus einem
anderen Arbeitsverzeichnis gestartet; die Quelle wird relativ zum
Prüfprogramm aufgelöst. Abhängigkeiten: NumPy, SymPy und SciPy.

Die symbolischen Identitäten und die vollständige W-Struktur werden
exakt geprüft. Zusätzlich läuft ein unabhängiger numerischer Vergleich
auf dem vollen 2076-dimensionalen N=2-Hamiltonoperator, direkt aus W
aufgebaut. Er bestätigt beide vollständigen Endzustände, die
Wahrscheinlichkeiten, Normerhaltung, verschwindende Bosonamplituden und
die Negativkontrolle. Der abschließende Lauf enthält \textbf{662 exakte
und 10 numerische Prüfbedingungen}. Prüfzahlen und Quellhash stehen in
\codeword{result.json}; \codeword{result\_OO.json} ist dessen
identischer optimierter Replay. Beide Ergebnisdateien haben SHA-256
\codeword{b37ae758117e4c35d1063e76756cd6a39f133c964d331a9d400e67a0b157a4a6}.

\textbf{Geschlossen:} Im ausdrücklich gewährten Instrumentenvertrag
existiert am selben ursprünglichen H ein kausaler
Moden-zu-Moden-Eingriff. Ein weiteres Hopping-Glied ist für diesen
N=2-Zeugen mathematisch nicht notwendig. Reine Paarumverteilung wird am
Endzeitpunkt getrennt von Speziesumwandlung bewiesen.

\textbf{Weiter offen:} Auswahl und Präparation dieses
N=2-Referenzzustands, ausführbare einzelne Modenphasen und
Besetzungsinstrumente aus dem nativen Operationssatz; Anschluss an den
N=64-Grundzustand und dessen geladenen Pol; Interpretation der
Modenalgebren als räumlich unabhängige Labore; Distanz, Raumdimension,
endliche Ausbreitungsgeschwindigkeit und gemeinsame Raumzeit. Die bloße
Tatsache, dass ein Kontrolloperator zur mathematischen Fock-Algebra
gehört, schließt diese Operationsfrage nicht.

\begin{center}\rule{0.5\linewidth}{0.5pt}\end{center}

\section{Teilbericht D: derselbe Dreizustandsraum für Schatten und
Eingriff}\label{teilbericht-d-derselbe-dreizustandsraum-fuxfcr-schatten-und-eingriff}

\section{Derselbe Dreizustandsraum trägt kausalen Eingriff und
Zustandstomografie}\label{derselbe-dreizustandsraum-truxe4gt-kausalen-eingriff-und-zustandstomografie}

\subsection{Exakte gemeinsame
Reduktion}\label{exakte-gemeinsame-reduktion}

Der ursprüngliche kausale Paarversuch bleibt vollständig im Raum

\[\adjustbox{max width=\linewidth}{$\displaystyle 
\mathcal K_3=\operatorname{span}\{|p_0\rangle,|R_7\rangle,|b_0\rangle\},
\quad |p_0\rangle=|4,57\rangle,\quad
|R_7\rangle=\frac1{\sqrt7}\sum_{q=1}^7s_q|p_q\rangle.
$}\]

Die acht Paare und ihre Vorzeichen stammen unverändert aus Zeile 0 des
eingefrorenen W; sie stehen im zugehörigen \codeword{REPORT.md}. Die
Isometrie J hat diese drei Vektoren als Spalten. Direkt aus W folgt

\[\adjustbox{max width=\linewidth}{$\displaystyle 
HJ=Jh_3,\qquad Z_4J=Jz_3,\qquad N_bJ=Jb_3,
$}\] \[\adjustbox{max width=\linewidth}{$\displaystyle 
h_3=\begin{pmatrix}0&0&-g\\0&0&\sqrt7g\\-g&\sqrt7g&\Delta\end{pmatrix},
\quad z_3=\operatorname{diag}(-1,1,1),\quad
b_3=\operatorname{diag}(0,0,1).
$}\]

Alle drei Intertwining-Identitäten sind exakt und haben keinen Restterm.
Somit haben beliebige Folgen aus freier H-Entwicklung und Z4 denselben
nativen dreidimensionalen Träger. Hier gilt immer N=2.

Die Empfängerbesetzung besitzt den exakten komprimierten \textbf{Effekt}

\[\adjustbox{max width=\linewidth}{$\displaystyle 
e_3=J^\dagger n_5J=\operatorname{diag}(0,1/7,0).
$}\]

Damit reproduziert der Dreizustandsraum ohne Näherung beide kausalen
Wahrscheinlichkeiten aus \codeword{REPORT.md}: x(4−x)/64 ohne mittleren
Impuls und 9x²/1024 mit Impuls. In beiden Armen sind am dort gewählten
Endzeitpunkt Nf=2 und Nb=0.

\subsection{Wesentliche Grenze: ein terminaler Effekt ist kein
geschlossenes
Instrument}\label{wesentliche-grenze-ein-terminaler-effekt-ist-kein-geschlossenes-instrument}

n5 selbst erhält diesen Dreizustandsraum nicht. Es gilt exakt

\[\adjustbox{max width=\linewidth}{$\displaystyle 
(n_5J-Je_3)^\dagger(n_5J-Je_3)=\operatorname{diag}(0,6/49,0).
$}\]

Eine tatsächliche nichtselektive projektive n5-Messung lässt vom Zustand
R7 das Gewicht \textbf{12/49} außerhalb von K3 zurück. Deshalb gilt die
gemeinsame Reduktion für Entwicklungen und Impulse \textbf{bis zur
abschließenden Messung}. Eine weitere Verwendung desselben Exemplars
nach diesem Instrument wird nicht behauptet. Die Tomografie verwendet
jeweils frische Exemplare desselben präparierten Zustands und genau eine
abschließende Besetzungsmessung.

Ein tatsächlich grobes Lüders-Instrument für den Projektor auf die
gesamten sieben anderen Paarblätter würde K3 erhalten; sein Effekt ist
in K3 diag(0,1,0). Das ist eine zusätzliche mögliche
Instrumentenanforderung. Ein feines Auslesen der sieben Blattlabels und
anschließendes Wegwerfen der Labels ist ein anderes Instrument: Es lässt
aus R7 das Gewicht 6/7 aus K3 heraustreten. Keine dieser
Implementierungen wird als nativ verfügbar vorausgesetzt, sofern sie
nicht ausdrücklich gewährt wird.

\subsection{Vollständige terminale Tomografie auf demselben
Träger}\label{vollstuxe4ndige-terminale-tomografie-auf-demselben-truxe4ger}

Definiere für hermitesche Operatoren

\[\adjustbox{max width=\linewidth}{$\displaystyle 
\mathcal L(O)=i[h_3,O],\qquad\mathcal C(O)=z_3Oz_3.
$}\]

Die folgende Liste besitzt neun linear unabhängige hermitesche
Operatoren:

\[\adjustbox{max width=\linewidth}{$\displaystyle 
\boxed{e_3,b_3,\mathcal Le_3,\mathcal Lb_3,
\mathcal L^2e_3,\mathcal L^2b_3,\mathcal L^3e_3,
\mathcal C\mathcal L^2e_3,\mathcal C\mathcal L^2b_3.}
$}\]

Schreibt man O als reellen Koordinatenvektor

\[\adjustbox{max width=\linewidth}{$\displaystyle 
(O_{00},O_{11},O_{22},\Re O_{01},\Re O_{02},\Re O_{12},
\Im O_{01},\Im O_{02},\Im O_{12}),
$}\]

hat die Matrix dieser neun Spalten die Determinante

\[\adjustbox{max width=\linewidth}{$\displaystyle 
\boxed{\det M=\frac{8\Delta^3g^{10}}{343}\ne0}
\qquad(\Delta g\ne0).
$}\]

Damit spannt der Orbit aus zeitentwickelten terminalen Effekten und
lokaler Impulskonjugation \textbf{Herm(3)} vollständig auf. Der bereits
bewiesene kausale Zeuge und diese Schattenrekonstruktion benutzen
denselben ursprünglichen Tensor, Hamiltonoperator, Träger, Phasenimpuls
und dieselben Besetzungsobservablen. Es handelt sich um
Zustandstomografie bei bekanntem H und bekannten Instrumenten, nicht um
eine gleichzeitige unbekannte Hamilton- und Instrumentenkalibrierung.

Die Operatorwörter sind durch Ableitungen von wirklichen
Endwahrscheinlichkeiten zugänglich:
\(\partial_t^k p_e(0)=\operatorname{tr}(\rho\mathcal L^ke_3)\). Für die
beiden letzten Wörter wird Z4 \textbf{vor} der freien Entwicklung
angewandt;
\(p_e^Z(t)=\operatorname{tr}(\rho z_3e^{ith_3}e_3e^{-ith_3}z_3)\). Die
zugehörige zweite Ableitung liefert das verlangte Wort. Es genügt also
der Zugriff auf die Kurven von n5 und Nb mit und ohne anfänglichen
Impuls, mit Ableitungen bis zur Ordnung drei. Die vollständige
Bestimmung dieser Ableitungen wird als Messdatenanforderung benannt; ein
optimiertes endliches Abtast- und Schätzverfahren wird hier nicht
behauptet.

\subsubsection{Der Impuls ist für normierte Zustandstomografie nicht
zwingend}\label{der-impuls-ist-fuxfcr-normierte-zustandstomografie-nicht-zwingend}

Die rein autonome Zeitfamilie aus e3 und b3 besitzt Dimension acht. Ihre
einzige fehlende hermitesche Richtung kann durch

\[\adjustbox{max width=\linewidth}{$\displaystyle 
Q=\begin{pmatrix}6\Delta&\sqrt7\Delta&-g\\
\sqrt7\Delta&0&\sqrt7g\\-g&\sqrt7g&0\end{pmatrix}
$}\]

dargestellt werden: \([h_3,Q]=0\), \(\operatorname{tr}(e_3Q)=
\operatorname{tr}(b_3Q)=0\), aber \(\operatorname{tr}Q=6\Delta\ne0\).
Für die Differenz zweier normierter Zustände ist diese Richtung daher
ausgeschlossen. Tatsächlich bilden

\[\adjustbox{max width=\linewidth}{$\displaystyle 
I,e_3,b_3,\mathcal Le_3,\mathcal Lb_3,\mathcal L^2e_3,
\mathcal L^2b_3,\mathcal L^3e_3,\mathcal L^4e_3
$}\]

ebenfalls eine Basis; ihre Determinante ist 6Δ³g¹⁰/343. Der bekannte
Wert Trρ=1 ersetzt somit eine weitere Messrichtung. Der Impuls erlaubt
hier eine konkrete Basis mit Ableitungen nur bis Ordnung drei; aus Rang
acht allein folgt kein Tomografiehindernis für normierte Zustände.

\subsection{Bedingte Fehlergrenze}\label{bedingte-fehlergrenze}

Verwende dimensionslose Zeit u=Δt und den Prüfpunkt g/Δ=1/20. Sei y der
Vektor der neun oben angegebenen Operatorerwartungswerte und δy sein
Fehler. Der exakte lineare Inversenrechner liefert mit der
Frobenius-Normschranke

\[\adjustbox{max width=\linewidth}{$\displaystyle 
\|\widehat\rho-\rho\|_{\mathrm{HS}}
\le\frac{\sqrt{9818033}}2\|\delta y\|_2,
$}\] \[\adjustbox{max width=\linewidth}{$\displaystyle 
\frac12\|\widehat\rho-\rho\|_1
\le\frac{\sqrt{29454099}}4\|\delta y\|_2.
$}\]

Sind alle neun Dateneinträge mit Fehler höchstens ε bekannt, gilt für
die zweite Schranke 3√29454099·ε/4. Diese konservative Schranke belegt
Stabilität für die ausdrücklich gegebenen Daten, aber auch eine mögliche
starke Fehlerverstärkung bei schwacher Kopplung. Sie enthält
\textbf{keine} kostenlose Gewinnung genauer Zeitableitungen aus endlich
vielen verrauschten Messungen. Messdauer, Anzahl frischer Exemplare,
Ableitungsschätzung und mögliche Verbesserungen der Messauswahl bleiben
eigene Ressourcenfragen.

